% !PRINTED 68
\weberpage{68}
%<<CONT
bestimmen, fassen wir die Abh\"angigkeit der Funktion $t$ auch
von $\omega_1, \omega_2$ ins Auge und bezeichnen sie, indem wir $g_1, g_2$ als
gegebene von $\omega_1, \omega_2$ unabh\"angige Zahlen betrachten, mit
$t(u, \omega_1, \omega_2)$. Es zeigt sich dann, da\ss, wenn $h$ ein willk\"urlicher
Faktor ist,
\[
t(h\,u,\ h\,\omega_1,\ h\,\omega_2)
\]
gleichfalls den Bedingungen (1) gen\"ugt, und da\ss\ demnach

\weq{2}{\frac{t(hu,\ h\omega_1,\ h\omega_2)}{t(u,\ \omega_1,\ \omega_2)} = f(\omega_1,\ \omega_2)}

von $u$ unabh\"angig ist.

Die Funktion $f(\omega_1, \omega_2)$ kann aber immer in die Form gesetzt
 werden:
\[
f(\omega_1,\ \omega_2) = \frac{\varphi(h\omega_1,\ h\omega_2)}{\varphi(\omega_1,\ \omega_2)};
\]
man hat nur, wenn $t(u)$ f\"ur $u = 0$ nicht verschwindet:
\[
\varphi(\omega_1,\ \omega_2) = t(0,\ \omega_1,\ \omega_2),
\]
und wenn $t(0) = 0$ ist, etwa
\[
\varphi(\omega_1,\ \omega_2) = \omega_1\, t'(0,\ \omega_1,\ \omega_2)
\]
zu setzen, wenn $t'$ die Derivierte von $t$ nach $u$ bezeichnet. Indem
man also jetzt

\weq{3}{\frac{t(u,\ \omega_1,\ \omega_2)}{\varphi(\omega_1,\ \omega_2)}}

wieder mit $t(u, \omega_1, \omega_2)$ bezeichnet, kann man den Bedingungen (1)
noch die hinzuf\"ugen, da\ss\ f\"ur ein beliebiges $h$

\weq{4}{t(u,\ \omega_1,\ \omega_2) = t(hu,\ h\omega_1\ h\omega_2).} % [sic] -- kein Komma zwischen $h\omega_1$ und $h\omega_2$ wie im Original gedruckt

Die Funktion $t$ h\"angt also jetzt nur noch von zwei Ver\"anderlichen, n\"amlich den Verh\"altnissen $u : \omega_1 : \omega_2$, ab. Wir setzen in (4)
\[
h = \frac{1}{\omega_1},\quad \frac{u}{\omega_1} = v,\quad \frac{\omega_2}{\omega_1} = \omega,
\]
und erhalten

\weq{5}{t(u,\ \omega_1,\ \omega_2) = t\left(\frac{u}{\omega_1},\ 1,\ \frac{\omega_2}{\omega_1}\right) = \vartheta(v, \omega),}

worin $\vartheta$ ein neues Funktionszeichen ist.

Hierdurch ist also eine neue Funktion $\vartheta(v)$ definiert von nur
zwei Variablen $v, \omega$, deren erste unbeschr\"ankt ver\"anderlich ist,
w\"ahrend $\omega$ nach der im \S\,16 gemachten Voraussetzung einen
\gesperrt{positiven imagin\"aren} Bestandteil hat; $v$ hei\ss t das \gesperrt{Argument},
$\omega$ der \gesperrt{Modul} der $\vartheta$-Funktion.
